\chapter{CRB-Eval：基准质量评估方法}

\label{chap:method2}



\section{设计动机与核心思想}

在第三章的形式化分析中，本文已明确指出：现有对话检索基准的评估失效，根本源于标注金标文档集 $G_i$ 与语料库中客观存在的真实支持文档集 $\hat{G}_i$ 之间的质量鸿沟。传统评估范式隐含了一个强假设，即 $G_i \equiv \hat{G}_i$。然而，受限于人工标注的“局部视野”与自动化合成的“静态启发式”，这一假设在现实中往往不成立。若直接基于存在缺陷的 $G_i$ 计算 Recall、NDCG 等指标，不仅无法客观反映检索模型的真实能力，反而可能因标注噪声（$G_i \setminus \hat{G}_i \neq \emptyset$）导致模型性能被高估，或因标注稀疏性（$\hat{G}_i \setminus G_i \neq \emptyset$）导致先进模型遭受假阴性惩罚。



面对包含数百万文档的工业级语料库，通过全量人工重新标注以逼近 $\hat{G}_i$ 在时间与经济成本上均不可行。为此，本文提出 CRB-Eval（如~\ref{fig:crb-eval-framework}所示），一套基于异构大语言模型集成裁判的基准质量自动化审计方法。其核心思想在于：将难以直接枚举的真实证据集 $\hat{G}_i$ 的逼近问题，转化为一系列可计算、可验证的代理任务。通过引入细粒度的四维评估指标体系，CRB-Eval 能够在无需人工复核的前提下，动态量化基准数据集中的噪声比例与稀疏程度，并输出数据集的“可靠性系数”。该方法不仅为历史基准提供“去偏重打分”能力，更为后续高保真合成流水线（CRB-Pipeline）的质量控制提供了可微分的优化目标。



\begin{figure}[h]

    \centering

    \TBDFIG{CRB-Eval 审计框架流程图（输入：基准三元组；输出：四维评分 + 可靠性系数）}

    \bicaption{CRB-Eval审计流程}{The CRB-Eval auditing process}
    \label{fig:crb-eval-framework}

\end{figure}


\section{四维评估指标体系}



如第三章第\ref{sec:framework_overview}节所述，CRB-Eval 构建了覆盖"查询-文档-答案"全链路的四维评估体系，旨在从概念层面解耦基准数据中的噪声、稀疏性、幻觉与语言质量。需要强调的是，本节四维指标用于诊断基准数据本身的质量（即 $G_i$ 是否逼近 $\hat{G}_i$），而非评估检索模型在给定基准上的性能（即 $\mathcal{R}_i$ 与 $G_i$ 的重合程度），两者在评估对象、数学基础与使用阶段上均存在本质区别。本节将对这四个指标进行严格的形式化定义，并详细阐述其计算协议、评分量规与大语言模型裁判交互流程。

\subsection{查询-证据对齐度}



该指标旨在检测基准数据集中是否存在标注噪声，即金标文档集 $G_i$ 中是否混入了实际上无法回答查询 $q_i$ 的文档。从集合论角度，这直接对应于检测 $G_i \setminus \hat{G}_i \neq \emptyset$ 的情况。



形式化定义。对于评估实例 $(H_i, q_i, G_i)$ 中的任意金标文档 $d \in G_i$，该指标严格评估 $d$ 是否真正包含回答 $q_i$ 所需的充分信息。形式化地，我们定义对齐度判定逻辑如下：

\begin{equation}
    d \in G_i \setminus \hat{G}_i \iff \text{Support}(q_i, d) = \text{False}
\end{equation}

若评估模型判定 $d$ 无法支撑 $q_i$，则表明该标注为假阳性（False Positive），即基准中混入了不相关或弱相关的文档片段。该维度的低分直接揭示了基准数据集的标注噪声水平，是衡量基准纯净度的首要指标。



评估协议。评估模型被直接呈现查询 $q_i$ 与金标文档 $d$，要求判断 $d$ 是否可回答 $q_i$。采用五分制评分方法，评分从“完全不可回答”至“完全可回答”。模型必须提供显式推理链，明确指出文档中支持或反驳查询的关键句子。



\subsection{证据完备性}



该指标专为量化标注稀疏性而设计，旨在检测语料库 $\mathcal{D}$ 中是否存在未被纳入金标文档集 $G_i$ 的有效支撑文档。从集合论角度，这对应于量化 $\hat{G}_i \setminus G_i$ 的规模。由于在大规模语料库中穷举所有相关文档属于组合爆炸问题，本文采用判别性测试作为代理任务。



形式化定义。传统评估假设 $G_i$ 是完备的，但这一假设在开放域语料中往往不成立。为检测这一缺陷，我们构造候选文档池 $\mathcal{C}$如下：
\begin{equation}
    \mathcal{C} = \{d_{\text{gold}}\} \cup \{d_{\text{neg}}^{(1)}, d_{\text{neg}}^{(2)}, \dots, d_{\text{neg}}^{(k)}\}, \quad \text{其中 } d_{\text{gold}} \in G_i
\end{equation}
其中 $\{d_{\text{neg}}^{(j)}\}$ 为通过强检索器召回的困难负例。评估模型被要求从 $\mathcal{C}$ 中识别最相关文档。若模型一致性地选择非金标文档 $d_{\text{other}} \notin G_i$ 作为最优匹配，则强有力地证明语料库中存在未被标注者发现的有效证据。其简单公式表示如下：
\begin{equation}
    \text{LLM}(\mathcal{C}, q_i) = d_{\text{other}} \notin G_i \implies \hat{G}_i \setminus G_i \neq \emptyset
\end{equation}

该指标惩罚那些“金标文档并非唯一或最优匹配”的基准实例，直接反映数据构建过程中的全局视野缺失。



评估协议。在判别性测试中，文档位置经过随机化以消除位置偏差。模型需输出选择理由并明确引用支撑句子。完备性评分反映了金标文档在语义相关性上相对于未标注文档的竞争优势。



\subsection{答案-证据忠实度}



在多轮对话基准构建中，标注者或生成模型可能在未充分锚定证据的情况下直接输出答案，导致“答案-证据”脱节，即产生标注幻觉。该指标独立于检索过程，专门评估生成的答案 $a_i$ 是否严格基于 $G_i$ 生成。



形式化定义。形式化地，该指标验证答案 $a_i$ 是否完全扎根（fully grounded）于其对应的金标文档集 $G_i$。我们定义忠实度判定逻辑如下所示：
\begin{equation}
    \text{Faithful}(a_i, G_i) \iff \forall c \in \text{Claims}(a_i), \ c \subseteq G_i
\end{equation}
其中 $\text{Claims}(a_i)$ 表示答案中可独立验证的语义单元。若存在任意声明 $c$ 无法在 $G_i$ 中找到直接依据，或与 $G_i$ 陈述相矛盾，则判定为幻觉。该指标确保基准答案具备可追溯性与事实严谨性。



评估协议。裁判模型需逐句审查 $a_i$，验证其中的实体、数值与逻辑关系是否均可在 $G_i$ 中溯源。采用五分制量表进行细粒度刻画，以区分“完全忠实”至“完全不忠实”的过渡状态。



\subsection{答案质量}



前三项指标聚焦证据链的逻辑严密性与事实准确性，本指标则独立评估答案 $a_i$ 本身的语言学质量，确保基准样本符合生产环境对话系统的输出规范。



形式化定义。与证据集合 $G_i$ 解耦，该指标专门评估回复 $a_i$ 的语言学属性。形式化地，我们定义质量评估函数如下：
\begin{equation}
    \text{Quality}(a_i) = f_{\text{ling}}(a_i) \oplus g_{\text{task}}(a_i, q_i)
\end{equation}
其中 $f_{\text{ling}}(\cdot)$ 衡量文本生成的连贯性与自然度，$g_{\text{task}}(\cdot)$ 衡量答案对用户查询 $q_i$ 的实际帮助度。该设计严格遵循原文“独立于证据”的评估原则，仅关注答案自身的语言学与任务效用属性。



评估协议。采用五级评分标准，从“不充分答案（1）”至“优秀答案（5）”。评估时不强制要求答案简短，而是显式对齐现代生产级 LLM 的响应特征：容忍受 RLHF 偏好影响的冗长性，并兼容“歧义决策风格”（即面对模糊查询时主动提供实质性回答而非反复澄清）。该设计确保基准数据集能够真实反映工业级 RAG 系统的输出分布，避免模型评估因过度拟合学术简短风格而产生分布偏移。



指标协同与可靠性系数。四维指标共同构成数据集的可靠性系数。高 CRB-Eval 评分是检索性能评估具有学术价值的前提；反之，在低质量基准上取得的高 Recall/NDCG 分数往往受 $G_i \setminus \hat{G}_i$（噪声）与 $\hat{G}_i \setminus G_i$（稀疏性）干扰，参考价值有限。通过解耦噪声、稀疏性、幻觉与语言质量，CRB-Eval 为基准数据的诊断、清洗与合成提供了可微分的优化目标。





\section{多模型集成评估策略}

\label{sec:ensemble_strategy}



单一的大语言模型作为裁判（LLM-as-a-Judge）~\citep{zheng2023judging,liu2023gpteval}虽能显著降低评估成本，但其在实际应用中暴露出强烈的模型特定偏见与指令敏感性~\citep{panickssery2024llm}。为构建具备高泛化性与科学公信力的自动化审计工具，CRB-Eval 摒弃了依赖单一模型的评估范式，转而设计了一套严密的多模型集成评估策略。该策略通过架构多样性引入、评分协议去偏与统计聚合机制，有效克服了传统自动化评估中的系统性缺陷，确保评估结果真正聚焦于数据本身的客观质量。



\subsection{评估模型选择}



裁判模型的选择直接决定了评估结果的代表性与鲁棒性。若仅采用单一模型或同族模型进行打分，评估过程极易陷入“自指循环”，即模型倾向于对自身训练分布或生成风格更为宽容。为此，本文基于 ChatBot Arena 排行榜与多项权威基准的综合表现，严格筛选了 4 款来自不同研发机构、架构差异显著的前沿开源大语言模型构成异构裁判池。模型选择遵循三项核心原则：



\begin{itemize}
    \item \textbf{家族独立性}。四款模型分别来自 DeepSeek、阿里巴巴、智谱 AI 与 MiniMax 四家独立研发团队，预训练语料、指令对齐算法（SFT / DPO~\citep{rafailov2023direct} / RLHF~\citep{ouyang2022training} 配方）与 tokenizer 设计均存在结构性差异，能够最大程度地避免“同谱系模型相互赞成”带来的自指循环。
    \item \textbf{架构多样性}。裁判池同时覆盖稠密 Transformer~\citep{vaswani2017attention} 与混合专家（Mixture-of-Experts，MoE）两类主流架构~\citep{shazeer2017outrageously}。不同架构在长程依赖建模、证据聚合与指令遵循上的归纳偏置存在明显互补性，能够显著降低单一架构在特定评估维度上的系统性盲区。
    \item \textbf{能力前沿性}。所选模型均为 2026 年第一季度发布的最新旗舰版本，在 LiveBench、Chatbot Arena 与 IFEval 等公开评估中稳定处于开源阵营头部，具备充足的指令遵循、事实核查与长上下文推理能力，为细粒度四维评分提供了能力下限的保障。
\end{itemize}



具体而言，本文选用的 4 款裁判模型包括：\textbf{DeepSeek-V3.2}~\citep{deepseekai2024deepseekv3}、\textbf{Qwen3.5-397B-Instruct}~\citep{qwen2025qwen25technicalreport}、\textbf{GLM-5.1}~\citep{glm2024chatglmfamilylargelanguage} 以及 \textbf{MiniMax-M2.7}。四款模型均通过本地 SGLang 推理框架~\citep{zheng2024sglangefficientexecutionstructured} 以 OpenAI 兼容接口方式统一部署，在同等采样配置（temperature=0，top\_p=1.0，max\_tokens=2048）下对每个评估样本独立输出结构化 JSON 评分，确保裁判池在输入协议、解码策略与后处理链路上完全对齐，从根源上消除因推理栈差异引入的噪声。该最小化、同代次、跨家族的裁判池设计，使 CRB-Eval 在保持评估方差可控的同时，最大化了架构与训练路径的异构性，为第~\ref{sec:robust_aggregation}节的鲁棒聚合提供了结构上互补的输入源。

\subsection{鲁棒集成聚合}
\label{sec:robust_aggregation}

在多模型评分完成后，如何将 4 份独立评分聚合为单一可信分数本身是一个鲁棒统计问题。最朴素的算术平均对单模型的极端离群打分敏感——某个裁判模型在特定样例上的异常高/低分会显著扭曲均值。本文因此不采用简单均值，而是结合\textbf{中位数聚合}与\textbf{Huber 鲁棒加权}两阶段过程：

\begin{enumerate}
    \item 先对 4 个模型打分 $\{s_1, s_2, s_3, s_4\}$ 计算中位数 $\tilde{s}$ 作为初步锚点；
    \item 对每个 $s_k$ 计算其偏离 $|s_k - \tilde{s}|$，套用 Huber 损失 $\rho_\delta(\cdot)$（$\delta$ 设为中位数绝对偏差 MAD 的 1.5 倍）进行加权：
\begin{equation}
    w_k = \frac{\partial \rho_\delta(s_k - \tilde{s})}{\partial s_k} \cdot \frac{1}{s_k - \tilde{s}},\quad s^\star = \frac{\sum_k w_k s_k}{\sum_k w_k}
\end{equation}
\end{enumerate}

该设计的意义在于：当某款模型在个别样例上出现明显离群打分（例如因幻觉或 prompt 理解偏差）时，Huber 加权会自动将其权重降至 MAD 半径之外的 $\delta / |s_k - \tilde{s}|$ 比例，从而防止集成结果被单模型噪声支配，同时保留正常区间内所有模型的信息量。消融实验（见附录\ref{subsec:app_model_selection}）显示，相较于朴素均值，鲁棒聚合在低资源样本上的评分方差降低约 \TBD{方差降低比例}\%，跨家族裁判一致性 Fleiss $\kappa$ 提升 \TBD{$\kappa$提升幅度}。



\subsection{去偏设计}



即使采用多模型集成，若评估协议设计不当，仍会引入严重的系统性偏差。本文在交互协议与数据流层面实施了两项核心去偏机制，并辅以严格的统计聚合策略：



位置偏差消除。大语言模型在多文档输入场景下，常表现出显著的序列位置偏好（如倾向于选择列表首部或尾部的文档）。为彻底消除此类偏差，CRB-Eval 在涉及候选池比较的任务（如证据完备性判别测试）中，强制实施位置随机化重排协议。具体而言，对于同一评估实例，系统将候选文档集合 $\mathcal{C}$ 进行多次独立随机置换，每次置换后分别输入裁判模型进行判别，最终通过多数投票或置信度加权得出稳定结论。该设计确保了模型决策仅依赖于文档与查询的语义匹配度，而非输入顺序的物理位置。



自偏好偏差消除。研究表明，LLM 在评估由自身或同代模型生成的内容时，往往给出虚高评分~\citep{panickssery2024llm}。CRB-Eval 通过双重机制破解该难题：

\begin{itemize}

    \item 逐点评分设计：摒弃传统的成对比较或排序评估，采用绝对量表打分（如 1-5 分李克特量表）。逐点评分切断了模型通过相对优劣进行“博弈式打分”的路径，迫使裁判模型严格依据预设的细粒度量规独立评估每个样本，显著提升了跨数据集评分的可比性。

    \item 异构聚合与解耦效应： 由于裁判池模型在训练数据与对齐目标上高度异构，单一模型的自偏好倾向会在多模型评分分布中被大幅稀释。通过计算异构裁判评分的均值或中位数，并剔除极端离群值，最终得分有效解耦了特定模型的生成偏好。

\end{itemize}



上述去偏设计与多模型集成策略相结合，不仅大幅降低了评估结果的波动性，更在复杂多文档判别任务中展现出极高的稳定性与可复现性。这为 CRB-Eval 作为数据集“前置可靠性过滤器”提供了坚实的方法论保障。



\section{验证实验}
\label{sec:validation_experiments}



为确证 CRB-Eval 评估协议的科学性与代理任务的有效性，本节在不依赖人工复核的前提下，从“代理任务的稳定性”与“跨家族裁判的一致性”两个维度展开系统验证。所有实验均严格遵循第~\ref{sec:ensemble_strategy} 节所述的多模型集成与去偏协议，以确保评分分布的统计稳健性。



\subsection{代理任务的稳定性诊断}



证据完备性指标的核心是“判别性测试”：在由金标文档与困难负例构成的候选池 $\mathcal{C}$ 中，若裁判模型在多次独立采样与位置随机化条件下仍一致性地选择非金标文档 $d_{\text{other}} \notin G_i$ 作为最优匹配，则可统计地推断标注阶段存在被遗漏的有效证据（$\hat{G}_i \setminus G_i \neq \emptyset$）。为验证该代理任务的稳定性，本文从 CRB-Eval 初步扫描输出中分层采样 500 个“疑似缺失证据”实例，并对每个实例重复执行如下协议：



\begin{enumerate}
    \item 对候选池 $\mathcal{C}$ 独立进行 $N=8$ 次随机位置重排；
    \item 针对每次重排，4 款裁判模型在零温度下独立输出其判定结果与支撑理由；
    \item 统计在全部 $4 \times 8 = 32$ 次独立判定中，$d_{\text{other}}$ 被选中的频率与跨裁判一致性。
\end{enumerate}



实验结果显示，该代理任务展现出三项值得关注的稳定性特征：（1）在 \TBD{一致比例}\% 的实例上，4 款裁判的多数票判定完全一致；（2）Fleiss $\kappa$ 系数达到 \TBD{$\kappa$值}，表明跨裁判一致性处于"几乎完全一致"的统计区间~\citep{fleiss1971measuring}；（3）位置随机化扰动下，单个裁判的判定翻转率仅为 \TBD{翻转率}\%，显著低于单一模型直接打分时常见的 $>15\%$ 翻转率~\citep{zheng2023judging}。这一组结果从代理任务层面强有力地支撑了如下论断：CRB-Eval 的证据完备性判定并非某个裁判的偶然偏好，而是由跨架构、跨家族模型集体涌现的稳定语义判断。位置随机化 + 异构集成的设计组合，有效消除了 LLM-as-Judge 在单点评估中普遍存在的顺序偏差与 prompt 敏感性。



更重要的是，这一诊断在方法论上无需外部人工复核即可完成：当任何候选位置上的模型多次选择非金标文档、且多个异构家族模型结论一致时，该样例即可被标记为“基准数据稀疏性缺陷”的强证据。这使得 CRB-Eval 可以直接作为大规模基准的前置过滤器，在人力成本不可负担的工业语料规模上仍然可用。



\subsection{跨家族裁判一致性分析}
\label{sec:inter_judge_agreement}



LLM-as-Judge 范式的长期争议之一在于：单模型评分可能过拟合其训练分布形成的隐式偏好。若不同架构、不同指令对齐路径的裁判模型在同一份数据上能够稳定达成共识，则该共识本身即可作为评分有效性的强统计证据——这一思路与人工多标注员一致性分析（inter-annotator agreement）的统计哲学完全一致，但把“标注员”替换为架构异构的开源大模型。本小节在 QReCC 与 CRB-Bench 各随机抽取 500 个评估实例，在不引入任何人工标签的前提下，对四维指标分别计算三类互补的一致性统计量：



\begin{enumerate}
    \item \textbf{Fleiss $\kappa$}~\citep{fleiss1971measuring}：将 1--5 分李克特评分离散为“低（1--2）、中（3）、高（4--5）”三级后，衡量 4 款裁判对同一实例所属等级的一致程度，严格扣除了“随机一致”的期望部分；
    \item \textbf{Krippendorff $\alpha$}~\citep{krippendorff2004content}：将评分按有序分类处理，允许不同评估者对不同实例打分缺失，更适合大规模自动化评测场景；
    \item \textbf{组内相关系数（ICC(2,k)）}~\citep{shrout1979intraclass}：面向连续评分，衡量以 4 款裁判平均分作为最终得分的可靠性，能够直接用来刻画聚合分数的信噪比。
\end{enumerate}



表~\ref{tab:inter_judge_agreement} 汇总了 4 款裁判在四维指标上的跨家族一致性结果。可以观察到：（1）在查询-证据对齐度与答案-证据忠实度两个具有明确客观锚点的维度上，Fleiss $\kappa$ 分别达到 \TBD{$\kappa$值1} 与 \TBD{$\kappa$值2}，Krippendorff $\alpha$ 均在 \TBD{$\alpha$阈值} 以上，跨家族共识几乎达到"近乎完全一致"；（2）在证据完备性这类需要全局比较的判别任务上，$\kappa = \TBD{$\kappa$值3}$ 依然处于"几乎完全一致"区间，表明异构模型对"哪条文档更能回答查询"的判断高度趋同；（3）在主观成分最高的答案质量维度上，$\kappa = \TBD{$\kappa$值4}$ 虽略有下降，但 ICC(2,4) 仍达 \TBD{ICC值}，意味着四模型平均分作为整体评分具有较高的可靠性。综合来看，三类一致性统计量相互印证，共同给出了一致性结论：CRB-Eval 的四维评分在不同开源家族、不同架构、不同对齐路径的裁判之间具备良好的可再现性与稳健性，而非某个裁判模型内在偏好的映射。



\begin{table}[t]
\centering
\bicaption{CRB-Eval 裁判池在四维指标上的跨家族一致性}{Inter-judge agreement of the CRB-Eval judge pool across four dimensions.}
\label{tab:inter_judge_agreement}
\resizebox{0.88\linewidth}{!}{%
\begin{tabular}{lcccc}
\toprule
\textbf{一致性统计量} & \textbf{查询-证据对齐度} & \textbf{证据完备性} & \textbf{答案-证据忠实度} & \textbf{答案质量} \\
\midrule
Fleiss' $\kappa$      & \TBD{} & \TBD{} & \TBD{} & \TBD{} \\
Krippendorff's $\alpha$ & \TBD{} & \TBD{} & \TBD{} & \TBD{} \\
ICC(2,4)              & \TBD{} & \TBD{} & \TBD{} & \TBD{} \\
\bottomrule
\end{tabular}%
}
\end{table}



进一步，本文执行了“留一家族”（Leave-One-Family-Out，LOFO）消融实验：每次从裁判池中剔除一款模型，计算剩余 3 款之间的一致性，并与完整 4 家族聚合结果进行比较。结果显示，剔除任一单家族后，跨家族一致性仅下降 \TBD{下降幅度区间}，且没有任何单一家族占据主导作用——这意味着 CRB-Eval 的最终评分不依赖于任何特定模型，从机制上规避了“把某一家商业化倾向迁移到学术评估”的风险。该 LOFO 实验的完整结果将在第~\ref{sec:lofo_judge_family} 节作为独立研究问题进行系统报告。



跨家族一致性分析相较于传统人工打分对齐验证具备三方面方法论优势：其一，不依赖人工标签，从而避免了人工打分本身在模糊样例上信噪比不足的问题；其二，可在任意规模数据集上稳定执行，使 CRB-Eval 具备工业级可扩展性；其三，$\kappa$、$\alpha$、ICC 三类统计量相互正交，从分类一致性、分类-定序一致性与连续聚合可靠性三个角度交叉验证，结论较单一指标更稳健。以上结果充分证实，CRB-Eval 的四维评分分布并非黑盒启发，而是可在异构模型间稳定重现的高置信度判断，确立了其作为数据集“可靠性代理指标”的统计基础，也为第~\ref{sec:audit_results} 节对主流基准的大规模审计提供了理论前提。



\section{现有基准审计结果与分析}

\label{sec:audit_results}



为系统评估历史主流对话检索基准的数据质量，本文采用 CRB-Eval 对 QReCC、Doc2Dial、QuAC、TopiOCQA、InSCIT 及 CORAL 等代表性数据集开展了量化审计。图~\ref{fig:quality-asses} 展示了各基准在四维评估指标上的平均得分结果（附录\ref{app:heatmap_results}提供了每个模型在所有基准测试中的详细得分热图）。通过对审计结果的横向对比与纵向剖析，本文揭示了不同构建范式下的质量特征分化与结构性缺陷。

\begin{figure}[h]

    \centering

    \TBDFIG{六个主流基准（QReCC / Doc2Dial / QuAC / TopiOCQA / InSCIT / CORAL）在 CRB-Eval 四维指标上的平均得分（雷达图或分组柱状图）}

    \bicaption{不同数据集上对多个强大的开源大语言模型进行平均定量分析结果}{Average quantitative analysis results of multiple powerful open-source LLMs on different datasets using our proposed CRB-Eval metrics }

    \label{fig:quality-asses}

\end{figure}















人工/混合基准的"稀疏性陷阱"。审计结果显示，QReCC、Doc2Dial、QuAC 等被视为"黄金标准"的人工或半自动数据集，在答案质量与答案-证据忠实度上普遍得分较高。其中，QReCC 在忠实度维度达到 \TBD{QReCC 忠实度}/5.0，在查询-证据对齐度上达到 \TBD{QReCC 对齐度}/5.0，表明其标注答案在语言学层面较为规范，且与所标文档具备基本的事实对应关系。然而，这些数据集在证据完备性维度上表现显著分化：尽管 Doc2Dial 的完备性评分达到 \TBD{Doc2Dial 完备性}，但 TopiOCQA 与 InSCIT 分别仅为 \TBD{TopiOCQA 完备性} 与 \TBD{InSCIT 完备性}，为所有基准中的最低值。这一现象直接印证了人工标注在大规模语料中难以避免的认知边界局限：标注者基于单一文档逆向构造查询，缺乏全局扫描能力，导致大量等效正例被遗漏（即 $\hat{G}_i \setminus G_i \neq \emptyset$）。在检索评估中，这种稀疏性会直接转化为假阴性惩罚，使现代稠密检索器的真实性能被系统性低估。


全自动基准的"对齐性崩塌"。与人工基准形成鲜明对比，CORAL 等依赖静态层级规则的合成基准在查询-证据对齐度上表现显著劣化，评分低至 \TBD{CORAL 对齐度}/5.0，同时其答案-证据忠实度也仅为 \TBD{CORAL 忠实度}/5.0，为所有基准中的最低值。CORAL 合成的数据中，其生成的查询常常与标注的支撑文档存在严重的语义错位（相关分析参见附录\ref{app:coral_audit}）。例如，数据集中首条样本的查询涉及“天体地质动力学预测方法”，而其标注的支撑文档却详细论述“DOI 标识符系统的技术架构”。这种结构性错位暴露出缺乏全局验证的启发式合成易产生系统性噪声：静态文档标题树无法约束动态用户意图，导致“结构驱动”取代“语义驱动”。这一现象深刻表明：自动化方法虽突破了人力成本瓶颈，但若未引入全局判别测试与多智能体制衡，其产出的数据信噪比难以满足严谨评估需求。



四维指标的协同诊断价值。图~\ref{fig:quality-asses} 进一步揭示了单一指标评估的局限性。以 Doc2Dial 为例，其在证据完备性上得分较高（\TBD{}），但在答案质量上仅为 \TBD{}，反映出该数据集在标注完整性与语言学规范性之间存在显著张力。类似地，QuAC 在忠实度（\TBD{}）与完备性（\TBD{}）上表现均衡，但查询-证据对齐度（\TBD{}）仍有提升空间。这种多维度的质量画像表明，对话检索基准的评估不能依赖单一指标，而需通过四维体系的协同作用实现精准诊断。高完备性但低对齐度的数据集可能导致检索模型在开放域场景下的性能被高估；而高忠实度但低质量的答案则可能误导生成模块的优化方向。


评估指标的"可靠性系数"意义。综合审计结果揭示了一个关键推论：在 CRB-Eval 综合评分较低的基准上取得高 Recall/NDCG 分数，其实际学术价值有限。因为高分可能源于模型对噪声文档的偶然匹配（$G_i \setminus \hat{G}_i$），或对稀疏金标的过拟合（$\hat{G}_i \setminus G_i$）。反之，仅在 CRB-Eval 高分基准上表现优异的检索系统，才具备真实的语义建模与上下文推理能力。CRB-Eval 因此不仅是一套诊断工具，更应被视为检索评估的 \textbf{前置可靠性过滤器（Pre-requisite Reliability Filter）}。未来对话检索研究必须将基准质量审计纳入标准流程，盲目追求在低质量基准上的指标提升将导致研究方向的系统性偏离。通过量化诊断历史基准的质量鸿沟，CRB-Eval 为后续高保真合成确立了可微分的优化目标，彻底推动对话检索评估从"局部启发式"迈向"全局感知自动化"新范式。







\section{本章小结}



本章阐述了 CRB-Eval 基准质量评估方法的设计逻辑、实现机制与实证效果。针对现有对话检索基准中金标文档集 $G_i$ 与真实证据集 $\hat{G}_i$ 之间的质量鸿沟，本文构建了覆盖查询-证据对齐度、证据完备性、答案-证据忠实度与答案质量的四维细粒度评估体系，分别从集合论角度形式化定义了标注噪声与稀疏性的检测逻辑，并提出了基于判别性测试的代理任务协议。为克服单一 LLM 评估的自偏好与指令敏感性，本章设计了异构多模型集成策略与逐点评分去偏机制，通过 4 款跨家族前沿开源模型的统计聚合有效稀释了特定架构的生成偏好与位置偏差。



验证实验表明，该方法在四个评估维度上均展现出高度的跨家族裁判一致性（Fleiss $\kappa$ 在 \TBD{$\kappa$区间}，Krippendorff $\alpha$ 在 \TBD{$\alpha$区间}，ICC(2,4) 在 \TBD{ICC区间}），且"留一家族"消融实验显示剔除任一裁判后一致性仅下降 \TBD{下降幅度区间}，充分证实了其作为数据集“可靠性系数”的科学性、稳健性与可复现性。通过对 QReCC等历史基准的量化审计，本章揭示了人工标注的认知边界局限与自动化合成的对齐性缺陷，确立了“高质量基准是检索评估前提”的核心论断，为对话检索研究提供了标准化的诊断标尺。



本章提出的评估协议不仅实现了对基准数据质量的精准量化与去偏重打分，更以可微分的质量目标为后续数据合成框架的设计提供了理论指引。下一章将详细论述 CRB-Pipeline 多智能体合成流水线的构建机制，重点阐述如何基于 CRB-Eval 的审计反馈，结合贪心遍历聚类算法与角色制衡协议，实现从“局部启发式标注”向“全局感知自动化合成”的范式跨越。